\documentclas[11pt]{scrartcl}
\usepackage[usenames,dvipsnames,svgnames]{xcolor}
\usepackage[shortlabels]{enumitem}
\usepackage[framemethod=TikZ]{mdframed}
\usepackage{amsmath,amssymb,amsthm}
\usepackage{epigraph}
\usepackage[colorlinks]{hyperref}
\usepackage{microtype}
\usepackage{mathtools}
\usepackage[headsepline]{scrlayer-scrpage}
\usepackage{thmtools}
\usepackage{listings}
\usepackage{derivative}
\renewcommand{\epigraphsize}{\scriptsize}
\renewcommand{\epigraphwidth}{60ex}
\addtolength{\textheight}{3.14cm}
\ihead{\footnotesize\textbf{BCX-entry}}
\ohead{\footnotesize Updated Sun 19 Oct 2025 21:28:42 UTC}
\providecommand{\clubs}[1]{$#1\clubsuit$}
\providecommand{\clubg}[1]{\bgroup\color{green!40!black}[$#1\clubsuit$]\egroup}

\providecommand{\ol}{\overline}
\providecommand{\eps}{\varepsilon}
\providecommand{\half}{\frac{1}{2}}
\providecommand{\dang}{\measuredangle} %% Directed angle
\providecommand{\CC}{\mathbb C}
\providecommand{\FF}{\mathbb F}
\providecommand{\NN}{\mathbb N}
\providecommand{\QQ}{\mathbb Q}
\providecommand{\RR}{\mathbb R}
\providecommand{\ZZ}{\mathbb Z}
\providecommand{\dg}{^\circ}
\providecommand{\ii}{\item}
\providecommand{\alert}{\textbf}
\providecommand{\opname}{\operatorname}
\providecommand{\ts}{\textsuperscript}
\DeclareMathOperator{\sign}{sign}
% hacks for arc
\providecommand{\tarc}{\mbox{\large$\frown$}}
\providecommand{\arc}[1]{\stackrel{\tarc}{#1}}
\reversemarginpar
\providecommand{\printpuid}[1]{\marginpar{\href{https://otis.evanchen.cc/arch/#1}{\ttfamily\footnotesize\color{green!40!black}#1}}}

\mdfdefinestyle{mdbluebox}{roundcorner=10pt,innerbottommargin=9pt,
    linecolor=blue,backgroundcolor=TealBlue!5,}
\declaretheoremstyle[headfont=\sffamily\bfseries\color{MidnightBlue},
    mdframed={style=mdbluebox},]{thmbluebox}
\mdfdefinestyle{mdredbox}{frametitlefont=\bfseries,innerbottommargin=8pt,
    nobreak=true,backgroundcolor=Salmon!5,linecolor=RawSienna,}
\declaretheoremstyle[headfont=\bfseries\color{RawSienna},
    mdframed={style=mdredbox},headpunct={\\[3pt]},postheadspace=0pt,]{thmredbox}
\mdfdefinestyle{mdgreenbox}{linecolor=ForestGreen,backgroundcolor=ForestGreen!5,
    linewidth=2pt,rightline=false,leftline=true,topline=false,bottomline=false,}
\declaretheoremstyle[headfont=\bfseries\sffamily\color{ForestGreen!70!black},
    mdframed={style=mdgreenbox},headpunct={ --- },]{thmgreenbox}
\mdfdefinestyle{mdblackbox}{linecolor=black,backgroundcolor=RedViolet!5!gray!5,
    linewidth=3pt,rightline=false,leftline=true,topline=false,bottomline=false,}
\declaretheoremstyle[mdframed={style=mdblackbox}]{thmblackbox}
\declaretheorem[style=thmredbox,name=Problem]{problem}
\declaretheorem[style=thmredbox,name=Required Problem,sibling=problem]{reqproblem}
\declaretheorem[style=thmbluebox,name=Theorem,numberwithin=problem]{theorem}
\declaretheorem[style=thmbluebox,name=Lemma,sibling=theorem]{lemma}
\declaretheorem[style=thmbluebox,name=Theorem,numbered=no]{theorem*}
\declaretheorem[style=thmbluebox,name=Lemma,numbered=no]{lemma*}
\declaretheorem[style=thmgreenbox,name=Claim,sibling=theorem]{claim}
\declaretheorem[style=thmgreenbox,name=Claim,numbered=no]{claim*}
\declaretheorem[style=thmblackbox,name=Remark,sibling=theorem]{remark}
\declaretheorem[style=thmblackbox,name=Remark,numbered=no]{remark*}
\declaretheorem[style=thmgreenbox,name=Definition,sibling=theorem]{definition}
\declaretheorem[style=thmgreenbox,name=Definition,numbered=no]{definition*}
\declaretheorem[style=thmblackbox,name=Example,sibling=theorem]{example}
\declaretheorem[style=thmblackbox,name=Example,numbered=no]{example*}

\newenvironment{walkthrough}{\noindent\textbf{\color{green!40!black}Walkthrough.}}{}
\newlist{walk}{enumerate}{3}
\setlist[walk]{label=\bfseries (\alph*)}
\newenvironment{subproof}[1][Proof]{%
\begin{proof}[#1] \renewcommand{\qedsymbol}{$\blacksquare$}}%
{\end{proof}}
\newenvironment{soln}{\begin{proof}[Solution]}{\end{proof}}

\usepackage{asymptote}
\begin{asydef}
size(8cm); // set a reasonable default
usepackage("amsmath");
usepackage("amssymb");
settings.tex="pdflatex";
settings.outformat="pdf";
import geometry;
void filldraw(picture pic = currentpicture, conic g, pen fillpen=defaultpen, pen drawpen=defaultpen) { filldraw(pic, (path) g, fillpen, drawpen); }
void fill(picture pic = currentpicture, conic g, pen p=defaultpen) { filldraw(pic, (path) g, p); }
pair foot(pair P, pair A, pair B) { return foot(triangle(A,B,P).VC); }
pair centroid(pair A, pair B, pair C) { return (A+B+C)/3; }
\end{asydef}

\newcommand{\goals}[2]{\bgroup
\sffamily\small \emph{Instructions}: Solve \clubg{#1}.
If you have time, solve \clubg{#2}.\egroup\par}
%% 426c616e6b204c615465587e
\begin{document}
\title{Submission for BCX-ENTRY}
\subtitle{OTIS (internal use)}
\author{Samuel Brandão}
\date{\today}
\maketitle

\begin{example*}[\href{https://aops.com/community/p3355365}{NIMO Winter 2014/2}, $0\clubsuit$]
	Determine, with proof, the smallest positive integer $c$
	such that for any positive integer $n$,
	the decimal representation of the number $c^n+2014$ has digits all less than $5$.
\end{example*} \printpuid{14NIMOW2}

\begin{walkthrough}
	This is quite easy; it's just meant to give a concrete example
	of a bi-directional problem.
	\begin{walk}
		\ii Figure out the answer.
		\ii Show that this choice of $c$ actually works.
		\ii Manually verify that none of the smaller $c$ work.
	\end{walk}
\end{walkthrough}
%% You're not expected to write up walkthroughs (unless you really want to).
%% The source is just for your reference.

\begin{example*}[HMMT 2016 T4, $0\clubsuit$]
	Let $n > 1$ be an odd integer.
	On an $n \times n$ chessboard the center square
	and four corners are deleted.  We wish to group
	the remaining $n^2-5$ squares into $\frac12(n^2-5)$ pairs,
	such that the two squares in each pair intersect at exactly one point
	(i.e.\ they are diagonally adjacent, sharing a single corner).

	For which odd integers $n > 1$ is this possible?
\end{example*} \printpuid{16HMMTT4}

\begin{walkthrough}
	Let's do some cases first.
	\begin{walk}
		\ii Can one do $n=3$?
		\ii Can one do $n=5$?
		\ii Can one do $n=7$?
	\end{walk}
	In fact, for most $n$ the task is impossible.
	This is a parity argument:
	we seek a coloring the cells by black and white (not the usual checkerboard)
	so that any valid pair has different colors.
	\begin{walk}[resume]
		\ii Find a coloring of the squares by black and white
		so that diagonally adjacent squares are opposite colors.
		(Optionally, find all such colorings.)
		\ii Use this to narrow down the set
		of possible $n$ to two values.
		\ii Wrap up the problem using your earlier work.
	\end{walk}
\end{walkthrough}
%% You're not expected to write up walkthroughs (unless you really want to).
%% The source is just for your reference.

\begin{example*}[\href{https://aops.com/community/p12189456}{JMO 2019/1}, $0\clubsuit$]
	There are $a+b$ bowls arranged in a row,
	numbered $1$ through $a+b$,
	where $a$ and $b$ are given positive integers.
	Initially, each of the first $a$ bowls contains an apple,
	and each of the last $b$ bowls contains a pear.
	A legal move consists of moving an apple from bowl $i$ to bowl $i+1$
	and a pear from bowl $j$ to bowl $j-1$,
	provided that the difference $i-j$ is even.
	We permit multiple fruits in the same bowl at the same time.
	The goal is to end up with the first $b$ bowls each containing a pear
	and the last $a$ bowls each containing an apple.
	Show that this is possible if and only if the product $ab$ is even.
\end{example*} \printpuid{19JMO1}

\begin{walkthrough}
	First we show that if $ab$ is even then the goal is possible.
	The basic idea is to use induction.
	\begin{walk}
		\ii If $\min(a,b) \ge 1$, and $a$ and $b$ are opposite parity,
		show that in one swap one can reduce from $(a,b)$ to $(a-1, b-1)$.
		\ii If $\min(a,b) \ge 2$, and $a$ and $b$ are both even,
		show that in two swaps one can reduce from $(a,b)$ to $(a-2, b-2)$.
		\ii Formulate a set of base cases and complete the proof via induction.
	\end{walk}
	Conversely, we need to show the task is impossible if $ab$ is odd.
	\begin{walk}[resume]
		\ii Let $X$ denote the number of apples in odd-numbered bowls,
		and let $Y$ denote the number of pears in odd-numbered bowls.
		Find a relation between $X$ and $Y$ that does not change under the operation.
		\ii Use this to show that the task is impossible when $ab$ is odd.
	\end{walk}
\end{walkthrough}
%% You're not expected to write up walkthroughs (unless you really want to).
%% The source is just for your reference.

\begin{example*}[\href{https://aops.com/community/p3160559}{Shortlist 2012 C1}, $0\clubsuit$]
	There are $n$ positive integers written in a row.
	Iteratively, Alice chooses two adjacent numbers $x$ and $y$
	such that $x>y$ and $x$ is to the left of $y$,
	and replaces the pair $(x,y)$ by either $(y+1,x)$ or $(x-1,x)$.
	Prove that she can perform at most $n^n$ such steps.
\end{example*} \printpuid{12SLC1}

\begin{walkthrough}
	A weaker result is easier to see:
	\begin{walk}
		\ii Look up the term ``lexicographic order'' if you don't know what it means.
		\ii Show the process must terminate, ignoring the bound of $n^n$.
	\end{walk}
	The hard part is to get the bound $n^n$,
	which doesn't depend on how big the numbers are,
	rather only depends on $n$ itself.

	The idea is that we need to pay attention to the \emph{relative order}
	of the numbers, rather than the numbers themselves.
	After all, the numbers on the board can be as large as Alice wants.

	So for each board state $B$,
	we define a permutation $\pi_B$ on $\{1,\dots,n\}$
	where the number in the $i$th position of $B$
	is the $\pi_B(i)$th smallest number. For example,
	\[ B = (1337, 42, 2012, 1000, 7)
		\longmapsto \pi_B = 42531. \]
	There is a little wrinkle: what do we do with ties?
	\begin{walk}[resume]
		\ii Choose a convention for breaking ties, so that each $B$ gives
		an unambiguous $\pi_B$, even if some numbers of $B$ are equal.
		\ii Under your convention, is $\pi_B$ ``monotonic'' in the lexicographic order?
		\ii If you answered ``yes'' to (d), then prove it.
		If you answered ``no'', give a different answer to (c) and try again.
		\ii Prove that the process terminates in at most $n!$ steps.
	\end{walk}
\end{walkthrough}
%% You're not expected to write up walkthroughs (unless you really want to).
%% The source is just for your reference.

\begin{example*}[\href{https://aops.com/community/p4774079}{USAMO 2015/4}, $0\clubsuit$]
	Steve is piling $m\geq 1$ indistinguishable stones
	on the squares of an $n\times n$ grid.
	Each square can have an arbitrarily high pile of stones.
	After he finished piling his stones in some manner,
	he can then perform \emph{stone moves}, defined as follows.
	Consider any four grid squares, which are corners of a rectangle,
	i.e.\ in positions $(i, k)$, $(i, l)$, $(j, k)$, $(j, l)$
	for some $1\leq i, j, k, l\leq n$, such that $i<j$ and $k<l$.
	A stone move consists of either removing one stone from each of
	$(i, k)$ and $(j, l)$ and moving them to $(i, l)$ and $(j, k)$ respectively,
	or removing one stone from each of $(i, l)$ and $(j, k)$
	and moving them to $(i, k)$ and $(j, l)$ respectively.

	Two ways of piling the stones are equivalent if they can be obtained
	from one another by a sequence of stone moves.
	How many different non-equivalent ways can Steve pile the stones on the grid?
\end{example*} \printpuid{15AMO4}

\begin{walkthrough}
	Despite its placement, this is not an especially quick problem,
	which is why it is saved as the last walkthrough.
	It is really rather easy to mess up some details,
	and moreover, it also takes a while to write-up depending on how you do it.

	On the other hand, the statement is pretty long-winded.
	Here's a tl;dr: stone moves look like the thing below,
	count the number of equivalence classes under stone moves.
	\begin{center}
		\begin{asy}
			size(5cm);
			for (int i=0; i<=6; ++i) {
					draw((0,i)--(6,i), gray);
					draw((i,0)--(i,6), gray);
				}
			draw((1.5,3.5)--(1.5,1.9), red+1.1, EndArrow(TeXHead));
			draw((4.5,1.5)--(4.5,3.1), red+1.1, EndArrow(TeXHead));
			draw(circle( (1.5, 1.5), 0.3), red+dashed);
			draw(circle( (4.5, 3.5), 0.3), red+dashed);
			filldraw(circle( (1.5, 3.5), 0.3), gray, black+1.4);
			filldraw(circle( (4.5, 1.5), 0.3), gray, black+1.4);
		\end{asy}
	\end{center}

	This walkthrough will present the cleanest approach I know of,
	due to Ankan Bhattacharya.
	But you should be aware that most solutions to the problem are much worse.

	\begin{walk}
		\ii Let $r_i$ denote the number of stones in the $i$th row,
		and $c_j$ the number of stones in the $j$th column.
		Show that $r_i$ and $c_j$ never change.

		\ii What are $\sum r_i$ and $\sum c_j$?

		\ii Show that the number of $2n$-tuples of integers
		$(r_1, \dots, r_n, c_1, \dots, c_n)$ satisfying
		\[ \sum r_i = \sum c_j = (\text{answer to (b)}) \]
		is $\binom{n+m-1}{m}^2$.
	\end{walk}
	So the classic mistake is to assume that (c) gives the answer.
	In truth, this is only the beginning of the problem.
	To see why, let's agree that the \emph{signature}
	of a piling is the tuple described in (c).

	Thus we have checked that if two pilings are equivalent,
	then they have the same signature.
	But this does not mean the number of piling methods
	is equal to the number of signatures!\footnote{For example,
		here is another invariant:
		the total number of stones is $m$, which takes on only one value.
		But that certainly does not mean the answer is $1$.}
	\begin{walk}[resume]
		\ii There are two more statements we have to prove
		to finish the problem. What are they?
	\end{walk}

	We'll actually now remodel the problem as follows.
	Forget about the entire grid.
	Instead, consider a blackboard where we write $(x,y)$
	for every stone in row $x$ and column $y$.
	Thus there should be exactly $m$ ordered pairs
	on the blackboard, one for each stone.

	Thus, a stone move amounts to switching the $y$-coordinates
	of two ordered pairs.

	\begin{walk}[resume]
		\ii Consider two pilings which have the same signature.
		Describe an algorithm to reach one from the other,
		using the blackboard analogy instead of the grid.

		\ii Moreover, show that every possible signature is achievable.
		(Why is this step necessary?)

		\ii Put everything together to complete the solution.
	\end{walk}
\end{walkthrough}
%% You're not expected to write up walkthroughs (unless you really want to).
%% The source is just for your reference.

\newpage

% ========================================
\section*{Practice problems}
\goals{42}{55}

\epigraph{It's not denial. I'm just selective about the reality I accept.}
{Calvin in \emph{Calvin and Hobbes}}
Warning: this unit is not short.
Tough coach makes the beginners sweat too :P
%%fakesubsection{When is it possible?}

\begin{problem}[\href{https://aops.com/community/p3355368}{NIMO Winter 2014/3}, $2\clubsuit$]
The numbers $1$, $2$, \dots, $10$ are written on a board.
Every minute, one can select three numbers
$a$, $b$, $c$ on the board, erase them,
and write $\sqrt{a^2+b^2+c^2}$ in their place.
This process continues until no more numbers can be erased.
What is the largest possible number
that can remain on the board at this point?
\end{problem} \printpuid{14NIMOW3}

%% Type your solution to NIMO Winter 2014/3 (\href{https://otis.evanchen.cc/arch/14NIMOW3/}{14NIMOW3}), proposed by Evan Chen here ...

%% --------------------------------------------------

\begin{problem}[\href{https://aops.com/community/p12752840}{Shortlist 2018 N1}, $5\clubsuit$]
Determine all pairs $(m, n)$ of positive integers
for which there exists a positive integer $s$
such that $sm$ and $sn$ have an equal number of divisors.
\end{problem} \printpuid{18SLN1}

%% Type your solution to Shortlist 2018 N1 (\href{https://otis.evanchen.cc/arch/18SLN1/}{18SLN1}), proposed by Ukraine here ...

%% --------------------------------------------------

\begin{reqproblem}[\href{https://aops.com/community/p2254765}{USAMO 2011/2}, $9\clubsuit$]
	An integer is assigned to each vertex of a regular pentagon
	so that the sum of the five integers is $2011$.
	A turn of a solitaire game consists of subtracting an integer $m$
	(not necessarily positive) from each of the integers at two neighboring vertices
	and adding $2m$ to the opposite vertex, which is not adjacent
	to either of the first two vertices.
	(The amount $m$ and the vertices chosen can vary from turn to turn.)
	The game is won at a certain vertex if, after some number of turns,
	that vertex has the number $2011$ and the other four vertices have the number $0$.
	Prove that for any choice of the initial integers,
	there is exactly one vertex at which the game can be won.
\end{reqproblem} \printpuid{11AMO2}

%% Type your solution to USAMO 2011/2 (\href{https://otis.evanchen.cc/arch/11AMO2/}{11AMO2}), proposed by Sam Vandervelde here ...

%% --------------------------------------------------

\begin{problem}[Soviet Union 1978, $3\clubsuit$]
A token is placed on an $n\times n$ chessboard.
Each of two players moves it in turn to an orthogonally neighbouring cell.
A player loses if they move the token to a cell that was already visited
(this includes the cell the token started with).
\begin{enumerate}[(a)]
	\ii If the token starts on a corner cell, for each $n$, who wins?
	\ii If the token starts adjacent to a corner cell, for each $n$, who wins?
\end{enumerate}
\end{problem} \printpuid{ZC1AD077}

%% Type your solution to Soviet Union 1978 (\href{https://otis.evanchen.cc/arch/ZC1AD077/}{ZC1AD077}) here ...

%% --------------------------------------------------

%%fakesubsection{Mins and maxes}

\begin{problem}[INMO 2020/3, $3\clubsuit$]
Consider a nonempty set $S \subseteq \{0, 1, \dots, 9\}$.
It turns out that every sufficiently large positive integer
can be written as the sum of two positive integers $a$ and $b$,
where both $a$ and $b$ have decimal digits only in $S$.
How small can $|S|$ be?
\end{problem} \printpuid{20INMO3}

%% Type your solution to INMO 2020/3 (\href{https://otis.evanchen.cc/arch/20INMO3/}{20INMO3}) here ...

%% --------------------------------------------------

\begin{problem}[\href{https://aops.com/community/p6220314}{JMO 2016/4}, $3\clubsuit$]
Find, with proof, the least integer $N$ such that
if any $2016$ elements are removed
from the set $\{1, 2, \dots, N\}$,
one can still find $2016$ distinct numbers
among the remaining elements with sum $N$.
\end{problem} \printpuid{16JMO4}

The least integer $N$ that satisfies the statement is
\[
	\sum^{4032}_{i=2017} i = 6049 \cdot 1008 = 6097392.
\]

Notice that if we form pairs of numbers from the set with equal sum, there will be
at least $3024$ such pairs. Even if at most $2016$ pairs are destroyed, there will
still remain at least $1008$ pairs with equal sum, i.e., $2016$ numbers. Each pair
consists of the $x$th number from the left and the $x$th number from the right.
This way, each pair has sum $6048+1$ as in $(1,6048), (2,6047), \dots, (3024,3025)$.
$6049 \cdot 1008 = 6097392$.

$N$ can't be less than $6,097,392$ because the least possible sum is
\[
	\sum^{2016}_{i=1} i=2033136.
\]
However, the first $2016$ numbers of the set can be removed, changing the least
possible sum to
\[
	\sum^{4032}_{i=2017} i = 6097392.
\]


%% Type your solution to JMO 2016/4 (\href{https://otis.evanchen.cc/arch/16JMO4/}{16JMO4}), proposed by Gregory Galperin here ...

%% --------------------------------------------------

\begin{problem}[\href{https://aops.com/community/p357975}{IMO 1974/4}, $5\clubsuit$]
Consider a sequence $a_1 < a_2 < \dots < a_p$ of $p$ positive integers.
We wish to partition an $8 \times 8$ chessboard into $p$ non-overlapping rectangles
(whose sides are aligned with that of the grid) such that the $i$\ts{th} rectangle
has exactly $a_i$ black squares and exactly $a_i$ white squares inside it.
Find the largest value of $p$ for which such a decomposition could occur.
Moreover, determine all sequences $(a_1, a_2, \dots, a_p)$ that achieve this maximum.
\end{problem} \printpuid{74IMO4}

%% Type your solution to IMO 1974/4 (\href{https://otis.evanchen.cc/arch/74IMO4/}{74IMO4}) here ...

%% --------------------------------------------------

\begin{problem}[\href{https://aops.com/community/p23517195}{USEMO 2021/1}, $5\clubsuit$]
Let $n$ be a positive integer and
consider an $n \times n$ grid of real numbers.
Determine the greatest possible number of cells $c$ in the grid such that
the entry in $c$ is both strictly \emph{greater} than the average of $c$'s column
and strictly \emph{less} than the average of $c$'s row.
\end{problem} \printpuid{21USEMO1}

%% Type your solution to USEMO 2021/1 (\href{https://otis.evanchen.cc/arch/21USEMO1/}{21USEMO1}), proposed by Holden Mui here ...

%% --------------------------------------------------

\begin{problem}[MOP 2011, $5\clubsuit$]
Seven points $A_1$, \dots, $A_7$ are marked on a circle $\gamma$.
We draw all $\binom 72 = 21$ line segments
of the form $\ol{A_i A_j}$ for $1 \le i < j \le 7$.
They meet at $n$ distinct points in the interior of $\gamma$
(not on the boundary).
How small can $n$ be, over all possible choices of $A_i$?
\end{problem} \printpuid{11MOPR21}

%% Type your solution to MOP 2011 (\href{https://otis.evanchen.cc/arch/11MOPR21/}{11MOPR21}) here ...

%% --------------------------------------------------

\begin{problem}[\href{https://aops.com/community/p4649033}{Tournament of Towns 2005}, $3\clubsuit$]
Originally, every square of $8 \times 8$ chessboard contains a rook.
One by one, rooks which attack an odd number of other rooks are removed.
(Rooks may not jump over other rooks.)
Find the maximal number of rooks that can be removed.
(A rook attacks another rook if they are on the same row or column
and there are no other rooks between them.)
\end{problem} \printpuid{05TTFSA3}

%% Type your solution to Tournament of Towns 2005 (\href{https://otis.evanchen.cc/arch/05TTFSA3/}{05TTFSA3}) here ...

%% --------------------------------------------------

\begin{problem}[\href{https://aops.com/community/p33793661}{OTIS Mock AIME II 2025/11, by Jiahe Liu}, $5\clubsuit$]
At an informatics competition each student earns a score in $\{0, 1, \dots, 100\}$
on each of six problems, and their total score is the sum of the six scores (out of $600$).
Given two students $A$ and $B$, we write $A \succ B$
if there are at least five problems on which $A$ scored strictly higher than $B$.

Compute the smallest integer $c$ such that the following statement is true:
for every integer $n \ge 2$,
given students $A_1$, \dots, $A_n$ satisfying $A_1 \succ A_2 \succ \dots \succ A_n$,
the total score of $A_n$ is always at most $c$ points more than the total score of $A_1$.
\end{problem} \printpuid{25OIMEII11}

%% Type your solution to OTIS Mock AIME II 2025/11, by Jiahe Liu (\href{https://otis.evanchen.cc/arch/25OIMEII11/}{25OIMEII11}), proposed by Jiahe Liu here ...

%% --------------------------------------------------

\begin{problem}[\href{https://aops.com/community/p32899297}{MP4G 2024/14}, $3\clubsuit$]
The set $S$ contains $2024$ elements.
What is the size of a largest collection $C$ of subsets of $S$
with the property that the intersection of every three subsets in $C$ is nonempty,
but the intersection of every four is empty?
\end{problem} \printpuid{24MP4G14}

%% Type your solution to MP4G 2024/14 (\href{https://otis.evanchen.cc/arch/24MP4G14/}{24MP4G14}) here ...

%% --------------------------------------------------

%%fakesubsection{Just show it's possible or impossible}

\begin{problem}[\href{https://aops.com/community/p9543994}{Slovenia TST 2018}, $3\clubsuit$]
Let $n \ge 2$ be an integer.
There are $n^2$ marbles, each of which is one of $n$ colors,
but not necessarily $n$ of each color.
Prove that they may be placed in $n$ boxes
with $n$ marbles in each box,
such that each box contains marbles of at most two colors.
\end{problem} \printpuid{SVNTST}

%% Type your solution to Slovenia TST 2018 (\href{https://otis.evanchen.cc/arch/SVNTST/}{SVNTST}) here ...

%% --------------------------------------------------

\begin{problem}[\href{https://aops.com/community/p6022569}{Kazakhstan 2016, added by Tilek Askerbekov}, $3\clubsuit$]
Let $n\geq 2$ be a positive integer.
Prove that all divisors of $n$ can be written as a sequence $d_1,\dots,d_k$ such that
for any $1\leq i<k$ one of $\frac{d_i}{d_{i+1}}$ and $\frac{d_{i+1}}{d_i}$ is a prime number.
\end{problem} \printpuid{16KAZ101}

%% Type your solution to Kazakhstan 2016, added by Tilek Askerbekov (\href{https://otis.evanchen.cc/arch/16KAZ101/}{16KAZ101}) here ...

%% --------------------------------------------------

\begin{problem}[\href{https://aops.com/community/p25627665}{Shortlist 2021 A1}, $2\clubsuit$]
Let $n \ge 1$ be an integer,
and let $A$ be a subset of $\{0,1,2,\dots,5^n\}$.
If $|A|=4n+2$, prove that there exist $a,b,c \in A$
such that $a < b < c$ and $c+2a > 3b$.
\end{problem} \printpuid{21SLA1}

%% Type your solution to Shortlist 2021 A1 (\href{https://otis.evanchen.cc/arch/21SLA1/}{21SLA1}) here ...

%% --------------------------------------------------

\begin{reqproblem}[\href{https://aops.com/community/p23322931}{PAGMO 2021/5}, $9\clubsuit$]
	Celeste has an unlimited amount of each type of $n$ types of candy,
	numbered type $1$, type $2$, \dots, type $n$.
	Initially she takes $m>0$ candy pieces and places them in a row on a table.
	Then, she chooses one of the following operations (if available) and executes it:
	\begin{itemize}
		\ii She eats a candy of type $k$, and in its position in the row she
		places one candy type $k-1$ followed by one candy type $k+1$, indices modulo $n$.
		\ii She chooses two consecutive candies which are the same type, and eats them.
	\end{itemize}
	Find all positive integers $n$ for which Celeste can
	leave the table empty for any value of $m$
	and any configuration of candies on the table.
\end{reqproblem} \printpuid{21PAGMO5}

%% Type your solution to PAGMO 2021/5 (\href{https://otis.evanchen.cc/arch/21PAGMO5/}{21PAGMO5}) here ...

%% --------------------------------------------------

\begin{problem}[\href{https://aops.com/community/p611}{IMO 1997/4}, $3\clubsuit$]
An $n \times n$ matrix whose entries come
from the set $S = \{1, 2, \dots , 2n - 1\}$
is called a \emph{silver} matrix if,
for each $i = 1, 2, \dots , n$,
the $i$-th row and the $i$-th column together
contain all elements of $S$. Show that:
\begin{enumerate}[(a)]
	\ii there is no silver matrix for $n = 1997$;
	\ii silver matrices exist for infinitely many values of $n$.
\end{enumerate}
\end{problem} \printpuid{97IMO4}

%% Type your solution to IMO 1997/4 (\href{https://otis.evanchen.cc/arch/97IMO4/}{97IMO4}) here ...

%% --------------------------------------------------

%%fakesubsection{Counting}

\begin{problem}[\href{https://aops.com/community/p1860909}{JMO 2010/1}, $2\clubsuit$]
Let $P(n)$ be the number of permutations $(a_1, \dots, a_n)$
of the numbers $(1, 2, \dots, n)$ for which $ka_k$ is a perfect square
for all $1 \leq k \leq n$.
Find with proof the smallest $n$ such that $P (n)$ is a multiple of $2010$.
\end{problem} \printpuid{10JMO1}

%% Type your solution to JMO 2010/1 (\href{https://otis.evanchen.cc/arch/10JMO1/}{10JMO1}), proposed by Andy Niedermier here ...

%% --------------------------------------------------

\begin{problem}[\href{https://aops.com/community/p28982442}{USAMTS 5/3/29}, $5\clubsuit$]
David has a deck of $2n$ distinct cards,
each with one of $n$ colors so that each color appears twice.
He deals the cards face-up on a table, one at a time.
If there are ever two cards of the same color
both face-up at once, he removes both of them.

Of the $(2n)!$ possible original orderings of the deck,
determine how many have the following property:
at every point, the table contains cards of at most two distinct colors.
\end{problem} \printpuid{USMT5329}

%% Type your solution to USAMTS 5/3/29 (\href{https://otis.evanchen.cc/arch/USMT5329/}{USMT5329}) here ...

%% --------------------------------------------------

\begin{reqproblem}[\href{https://aops.com/community/p12071600}{Canada 2019/3}, $5\clubsuit$]
	Let $m$ and $n$ be positive integers.
	A $2m \times 2n$ grid of squares is colored
	in the usual chessboard fashion.
	Determine the number of ways to place $mn$
	counters on the \emph{white} squares,
	at most one counter per square,
	so that no two counters are diagonally adjacent.
\end{reqproblem} \printpuid{19CAN3}

%% Type your solution to Canada 2019/3 (\href{https://otis.evanchen.cc/arch/19CAN3/}{19CAN3}) here ...

%% --------------------------------------------------

%%fakesubsection{Extra}

\begin{problem}[\href{https://aops.com/community/p10191585}{EGMO 2018/4}, $5\clubsuit$]
Let $n \ge 3$ be an integer.
Several non-overlapping dominoes are placed on an $n \times n$ board.
The \emph{value} of a row or column is the number of dominoes
that cover at least one cell of that row or column.
A domino configuration is called \emph{balanced} if there exists
some $k \ge 1$ such that every row and column has value $k$.

Prove that a balanced configuration exists for every $n \ge 3$
and find the minimum number of dominoes needed
in such a configuration.
\end{problem} \printpuid{18EGMO4}

%% Type your solution to EGMO 2018/4 (\href{https://otis.evanchen.cc/arch/18EGMO4/}{18EGMO4}), proposed by Merlijn Staps (NLD) here ...

%% --------------------------------------------------

\begin{reqproblem}[\href{https://aops.com/community/p21498566}{JMO 2021/4}, $5\clubsuit$]
	Carina has three pins, labeled $A$, $B$, and $C$, respectively,
	located at the origin of the coordinate plane.
	In a \emph{move}, Carina may move a pin to
	an adjacent lattice point at distance $1$ away.
	What is the least number of moves that Carina can make
	in order for triangle $ABC$ to have area $2021$?
\end{reqproblem} \printpuid{21JMO4}

%% Type your solution to JMO 2021/4 (\href{https://otis.evanchen.cc/arch/21JMO4/}{21JMO4}), proposed by Brandon Wang here ...

%% --------------------------------------------------

\begin{problem}[\href{https://aops.com/community/p27349487}{USAMO 2023/5}, $3\clubsuit$]
Let $n\geq3$ be an integer. We say that an arrangement of the numbers
$1$, $2$, \dots, $n^2$ in an $n \times n$ table is \emph{row-valid}
if the numbers in each row can be permuted to form an arithmetic progression,
and \emph{column-valid} if the numbers in each column
can be permuted to form an arithmetic progression.

For what values of $n$ is it possible to transform any row-valid arrangement
into a column-valid arrangement by permuting the numbers in each row?
\end{problem} \printpuid{23AMO5}

%% Type your solution to USAMO 2023/5 (\href{https://otis.evanchen.cc/arch/23AMO5/}{23AMO5}), proposed by Ankan Bhattacharya here ...

%% --------------------------------------------------

\begin{reqproblem}[\href{https://aops.com/community/p1116177}{USAMO 2008/4}, $5\clubsuit$]
	For which integers $n \ge 3$ can one find a triangulation
	of regular $n$-gon consisting only of isosceles triangles?
\end{reqproblem} \printpuid{08AMO4}

%% Type your solution to USAMO 2008/4 (\href{https://otis.evanchen.cc/arch/08AMO4/}{08AMO4}), proposed by Gregory Galperin here ...

%% --------------------------------------------------

(Any set of $n - 3$ diagonals of $\mathcal{P}$ that
do not intersect in the interior of the polygon
determine a \textit{triangulation} of $\mathcal{P}$ into $n - 2$ triangles.)

\begin{problem}[\href{https://aops.com/community/p33958997}{XOOK 2025/1}, $9\clubsuit$]
For the Exeter disco party,
Oron wants to tile a $2024 \times 2024$ square in the middle of the dance floor
with $1 \times 1$ tiles which have arrows pointing either up, left, right, or down.
After placing down all the tiles in some configuration,
he realizes that he wants to rotate some of the tiles so that
no matter which tile a dancer starts on,
after following the arrows they end up leaving the grid through some fixed cell.
What's the minimum number of tiles he needs to rotate so that this is always possible?
\end{problem} \printpuid{25XOOK1}

%% Type your solution to XOOK 2025/1 (\href{https://otis.evanchen.cc/arch/25XOOK1/}{25XOOK1}), proposed by Royce Yao here ...

%% --------------------------------------------------

\end{document}
